\section{Conclusion}
We have in this thesis primarily focused on the $\alpha_1$ band decays of the \ce{^12C} \SI{17.76}{MeV} resonance. We also applied the main parts of our analysis to the \SI{16.62}{MeV} and \SI{18.35}{MeV} resonances; the former to validate our tools and replicate the results of Kuhlwein \cite{Kuhlwein}, and the latter since we had the data on hand anyways. \\

The results from both the Dalitz-plot and angular correlation analysis of the \SI{17.76}{MeV} resonance are in agreement: there is no trace of any $0^+$ state at this energy. Instead, a mix of $2^-$ and $3^-$ resonances was found to best describe our data. \\

Our results do not agree with the $0^+$ assignment of the \SI{17.76}{MeV} resonance. Based on a visual comparison, a correlation angle analysis, and finally a fit to its Dalitz plot, a $2^-$ assignation would be a much better match for its decay to the first excited state of \ce{^8Be}. In particular, \textit{no} trace of a $0^+$ state could be found. Our analysis does not extend to its ground state decays, but based on e.g. \cite{Segel}\cite{Munch}, a $0^+$ natural parity assignation is inevitable. A new $2^-$ resonance with an energy similar to \SI{17.76}{MeV} could explain these contradictory results. \\

Our results for the $2^-$ resonance agrees with \cite{Kuhlwein}, as long as the Coulomb interaction between the final $\alpha$ particles is neglected. The analysis of the $3^-$ resonance is inconclusive, since our simulations cannot possibly fully replicate the data. It is believed that this issue can be solved by fine-tuning the R-matrix parameters. In particular it is worth noting that our analysis does not disagree with the $2^-$ and $3^-$ assignations of the \SI{16.62}{MeV} and \SI{18.35}{MeV} resonances. 